% xcodebuild-simctl-cli.tex — Xcode from the command line for CI, as mechanisms: one invocation's
% anatomy, archive → export → upload, build once / shard / merge, asking the tools before guessing,
% parallel clones, the two repetition loops, prefix-stripped environments, reading the .xcresult and
% the exit status, who an export method lets install, headless signing, SwiftPM on CI, what a
% simulator really is, which Xcode answers, and Xcode Cloud's three scripts.
% Sources (checked 2026-10-11) — first checked 2026-10-10:
%   keith.github.io/xcode-man-pages (mirror of the Xcode 27.2 beta 2 man pages): xcodebuild.1 (SYNOPSIS:
%     -workspace needs -scheme, the rest bracketed; actions; -list, -showdestinations,
%     -showBuildSettings, -showTestPlans "Does not initiate a build"; destination keys, id XOR name,
%     OS = latest by default, generic/ prefix, several -destination = whole suite on each
%     concurrently, -parallelize-tests-among-destinations, -parallel-testing-worker-count;
%     -testProductsPath / -xctestrun refuse -workspace / -project; -only-testing Target[/Class[/method]]
%     and @response-file, -only-testing takes precedence over -skip-testing; -enumerate-tests,
%     -test-enumeration-style flat "can subsequently be passed to -skip-testing and -only-testing";
%     -test-iterations "the maximum number of iterations" with either loop, -retry-tests-on-failure
%     (max 3 assumed), -run-tests-until-failure (max 100 assumed), -test-repetition-relaunch-enabled
%     (a new process per repetition); -resultBundlePath "If the path already exists, xcodebuild will
%     exit with an error"; -exportArchive takes no action; -allowProvisioningUpdates (automatic:
%     creates and updates profiles, app IDs and certificates; manual: downloads profiles; needs an
%     account in Xcode's Accounts pane) and -authenticationKeyPath / -authenticationKeyID /
%     -authenticationKeyIssuerID; -xcconfig "override all other settings, including settings passed
%     individually on the command line", XCODE_XCCONFIG_FILE overrides those "and those in the file
%     passed with the -xcconfig option"; TEST_RUNNER_ prefix "passed, with its prefix stripped, to all
%     test runner processes"; "exits with codes defined by sysexits(3) … EX_SOFTWARE if the commands
%     given to xcodebuild fail"; -runFirstLaunch) · xcodebuild.xctestrun.5 (build-for-testing writes
%     the scheme's parameters to it; editable; TestBundlePath, TestHostPath, EnvironmentVariables,
%     CommandLineArguments) · xcresulttool.1 (get test-results summary / tests, export attachments,
%     merge) · xcode-select.1 (-s needs sudo, affects all users; DEVELOPER_DIR overrides it; the tool
%     shims invoke the selected directory's tool) · xcrun.1 (--find) · altool.1 (--upload-package)
%   github.com/fastlane/fastlane/issues/13190 (a CI log: "** TEST FAILED **" then "Exit status: 65")
%     · developer.apple.com/forums/thread/121458 (accepted answer quotes the man page: 70 for failed
%     commands; asker sees 65) · github.com/apache/cordova-ios/issues/966 ("xcodebuild: Command
%     failed with exit code 65") · discuss.bitrise.io/t/…/13664 ("** BUILD FAILED **" …
%     "xcodebuild: Command failed with exit code 65") · github.com/fastlane/fastlane/issues/29553
%     ("Unable to find a destination matching the provided destination specifier" … "Exit status: 70")
%   developer.apple.com/forums/thread/690647 (Xcode 13 notes: build numbers synced "in a staged copy",
%     "doesn't modify … your archive")
%   developer.apple.com/documentation/xcode-release-notes: xcode-13 (test repetition modes;
%     xcodebuild authenticates with App Store Connect API keys for automatic signing in headless
%     environments) · xcode-15 (-enumerate-tests; -exportArchive uploads with ASC keys) · xcode-16
%     (an unmatched -destination is no longer silently replaced; -downloadPlatform -exportPath,
%     -importPlatform; devicectl --device by name / UDID; manual export needs provisioningProfiles)
%   developer.apple.com/documentation/xcode: downloading-and-installing-additional-xcode-components
%     (-downloadPlatform -exportPath once, -importPlatform after -runFirstLaunch) ·
%     building-swift-packages-or-apps-that-use-them-in-continuous-integration-workflows (commit
%     Package.resolved; -disableAutomaticPackageResolution uses exactly it; SSH URLs + known_hosts;
%     -scmProvider system = the Mac's git and configuration) · writing-custom-build-scripts (ci_scripts
%     beside the project, the three names, post-xcodebuild runs "even if the xcodebuild command fails",
%     no sudo, secrets as asterisks) · distributing-your-app-for-beta-testing-and-releases (upload /
%     export, manage build number) · building-your-app-to-include-debugging-information (archive keeps
%     dSYMs, build UUID)
%   developer.apple.com/help/account: devices/devices-overview (100 devices per product family per
%     membership year; disabling one "won't increase your number of available devices"; restored at
%     the new membership year) · provisioning-profiles/create-an-ad-hoc-provisioning-profile (select
%     the registered devices the profile includes)
%   developer.apple.com/videos: wwdc2019/418 ("the iOS, watchOS or tvOS userspace but running on the
%     Mac kernel", own launchd, same filesystem, separate home; "memory and CPU limits … are not
%     simulated", a terabyte Mac → a terabyte simulator; list --json; prefer full identifiers; clone
%     after shutdown, APFS file cloning) · wwdc2020/10647 (privacy grant, push, io, status_bar
%     override, keychain add-root-cert, booted alias)
%   asciiwwdc.com/2018/sessions/403 (WWDC18 "What's New in Testing": the selected simulator is cloned,
%     "the original simulator is not used during testing … a sort of template"; one runner per clone;
%     distributed by class)
%   github.com/ios-control/ios-sim/issues/161 (simctl launch help: "set them in the calling
%     environment with a SIMCTL_CHILD_ prefix") · nshipster.com/simctl (xcrun simctl list --json;
%     list devices)
%   developer.apple.com/forums: 773749 (DTS: export method list, default debugging) · 763050 (DTS:
%     xcresulttool help prints the JSON schema) · 11313 (simctl erase) · 812475 (simctl openurl) ·
%     759317, 725786 (ExportOptions keys) · 739347 (-skipMacroValidation help text) · 694058, 699941
%     (-clonedSourcePackagesDirPath; default DerivedData/SourcePackages)
%   gnu.org/software/bash/manual Pipelines (status of the last command unless pipefail) ·
%     github.com/cpisciotta/xcbeautify README · github.com/XcodesOrg/xcodes README (.xcode-version)
% @labels: area=cicd kind=tooling platform=ios topic=build,testing,tooling discussion=300 review=documentation-like,low-quality,boring
% @tags: xcodebuild, simctl, devicectl, xcarchive, ipa, export-options-plist, xcresult, xcresulttool, xccov, test-sharding, parallel-testing, test-repetition, xctestplan, swift-package-manager, package-resolved, app-store-connect-api-key, xcode-select, xcode-cloud, ci-cd
\documentclass[8pt]{extarticle}
\usepackage{printup-sheet}
\usepackage{array}
\usetikzlibrary{shadows}

% Linked text: the website links it, paper prints the text only (paper holds no links).
\newcommand\link[2]{\pdfonly{#1}\htmlonly{\href{#2}{#1}}}

\newcommand\ct[1]{\texttt{#1}}
\newcommand\dd{\pdfonly{-\kern0pt-}\htmlonly{-{}-}}   % "--" without the ligature on paper, plain "--" on the web
\newcommand\hd[1]{\par\vspace{3pt}\noindent{\bfseries\color{sheetBlue}#1}\par\vspace{1pt}}

\tikzset{
  act/.style={box, draw=sheetOrange, fill=sheetOrange!10, font=\tiny\ttfamily, inner sep=1.5pt, align=center},
  fil/.style={draw=sheetBlue, thick, fill=white, font=\tiny, inner sep=1.5pt, align=left, rounded corners=1pt},
  lbl/.style={font=\tiny, text=black!75, inner sep=1pt, align=center},
  ttl/.style={font=\bfseries\scriptsize, text=sheetBlue, anchor=west},
  cmd/.style={font=\ttfamily\tiny, inner sep=0.6pt, anchor=base west},
  brc/.style={decorate, decoration={brace, mirror, amplitude=2.5pt}, draw=sheetGrey},
  ok/.style={circle, fill=sheetGreen, inner sep=0pt, minimum size=2mm},
  ko/.style={circle, fill=sheetRed, inner sep=0pt, minimum size=2mm},
  st/.style={draw=sheetBlue, thick, rounded corners=2pt, fill=sheetBlue!6, font=\tiny, inner sep=1.5pt, align=center},
}

\begin{document}
\pdfonly{\folioslug{cicd/xcodebuild-simctl-cli}}

\sheettitle{xcodebuild, simctl \& friends — Xcode from the command line}{cicd · folio}

\oneliner{\ct{xcodebuild} runs one \textbf{action} of a \textbf{scheme} against a \textbf{destination}
and leaves files you can open — an \ct{.xcarchive}, an \ct{.ipa}, an \ct{.xcresult}. \ct{simctl}
drives simulators, \ct{devicectl} phones; \ct{xcode-select} / \ct{DEVELOPER\_DIR} pick \textbf{which
Xcode} answers.}

\vspace{2pt}
\noindent\begin{tikzpicture}[sheet]
  % ---- anatomy of one invocation ----
  % One command on two rows, as a shell line continuation. Every node is placed AT the previous
  % node's base east: `right=of` resets the anchor to west and walks the baseline down.
  \node[ttl] at (-0.15,5.5) {One invocation — a workspace needs \ct{-scheme}; the rest is optional};
  \node[cmd] (c1) at (-0.12,5.05) {xcodebuild};
  \node[cmd] (c2) at ([xshift=0.6mm]c1.base east) {-workspace A.xcworkspace};
  \node[cmd] (c3) at ([xshift=0.6mm]c2.base east) {-scheme App};
  \node[cmd] (c4) at ([xshift=0.6mm]c3.base east) {-configuration Debug};
  \node[cmd] (c6) at ([xshift=0.6mm]c4.base east) {-resultBundlePath r.xcresult};
  \node[cmd] at ([xshift=0.6mm]c6.base east) {\textbackslash};
  \path (c2.base west) ++(0,-0.83) coordinate (r2);
  \node[cmd] (c5) at (r2) {-destination 'platform=iOS Simulator,name=iPhone 17,OS=26.4'};
  \node[cmd, text=sheetOrange] (c7) at ([xshift=0.6mm]c5.base east) {\textbf{test}};
  \node[cmd] (c8) at ([xshift=0.6mm]c7.base east) {K=V};
  \draw[brc] (c2.south west) -- (c2.south east) node[midway, below=2.5pt, lbl] {the container\\(or \ct{-project})};
  \draw[brc] (c3.south west) -- (c3.south east) node[midway, below=2.5pt, lbl] {which targets,\\which test plan};
  \draw[brc] (c4.south west) -- (c4.south east) node[midway, below=2.5pt, lbl] {Debug or\\Release settings};
  \draw[brc] (c6.south west) -- (c6.south east) node[midway, below=2.5pt, lbl] {the results; the path\\must not exist yet};
  \draw[brc] (c5.south west) -- (c5.south east) node[midway, below=2.5pt, lbl] {where it runs: \texttt{platform}, then \texttt{id} or \texttt{name}; no \texttt{OS} means the latest;\\\texttt{generic/platform=iOS} = any device, none attached (archives)};
  \draw[brc] (c7.south west) -- (c7.south east) node[midway, below=2.5pt, lbl] {the\\action};
  \node[lbl, anchor=west, align=left] (kv) at ([xshift=4mm]c8.east) {a build setting for this run; \ct{-xcconfig f} overrides it,\\and \ct{XCODE\_XCCONFIG\_FILE} overrides both};
  \draw[flow, thin] (kv.west) -- (c8.east);
  % ---- release path ----
  \node[ttl] at (-0.15,3.42) {Release — each arrow a command, each box a file you can open};
  \node[fil, text width=16mm, minimum height=10mm] (src) at (0.72,2.78) {\textbf{source}\\\ct{.xcworkspace}\\scheme\\\ct{Package.resolved}};
  \node[act, text width=19.5mm, minimum height=10mm] (a1) at (2.98,2.78) {\textbf{archive}\\-archivePath A\\-destination\\generic/platform=iOS};
  \node[fil, text width=22mm, minimum height=10mm] (xa) at (5.5,2.78) {\textbf{\ct{A.xcarchive}}\\the \ct{.app} + every dSYM\\(matched by build UUID)\\keep one per shipped build};
  \node[act, text width=20mm, minimum height=10mm] (a2) at (8.2,2.78) {\textbf{-exportArchive}\\-archivePath A\\-exportPath out\\-exportOptionsPlist \textcolor{sheetRed}{E}};
  \node[fil, text width=21mm, minimum height=10mm] (ipa) at (10.95,2.78) {\textbf{\ct{out/App.ipa}}\\a zip: \ct{Payload/App.app}\\repackaged + signed\\for \ct{method}};
  \node[box, draw=sheetGreen!70!black, fill=sheetGreen!10, font=\tiny, text width=16mm, minimum height=10mm] (asc) at (14.05,2.78) {\textbf{App Store Connect}\\processing\\$\rightarrow$ TestFlight\\$\rightarrow$ review};
  \draw[flow] (src) -- (a1); \draw[flow] (a1) -- (xa); \draw[flow] (xa) -- (a2); \draw[flow] (a2) -- (ipa);
  \draw[hot] (ipa) -- node[lbl, above]{upload} (asc);
  \node[lbl, anchor=north west] at (2.0,2.2) {\textcolor{sheetRed}{E} = ExportOptions.plist; with \ct{destination=upload} the export step uploads instead of writing \ct{out/}};
  \node[lbl, anchor=north east] at (14.9,2.2) {or \ct{altool \dd upload-package}};
  % ---- test path: build once, shard, merge ----
  \node[ttl] at (-0.15,1.75) {Tests — build once, enumerate, shard across Macs, merge the results};
  \node[act, text width=17mm, minimum height=11mm] (b1) at (0.8,1.05) {\textbf{build-for-}\\\textbf{testing}\\-testProductsPath\\T.xctestproducts};
  \node[act, text width=21mm, minimum height=11mm] (en) at (3.35,1.05) {\textbf{test-without-building}\\-testProductsPath T\\-enumerate-tests\\-test-enumeration-\\style flat};
  \node[fil, text width=12mm, minimum height=11mm] (sf) at (5.65,1.05) {\ct{s1.txt}\\\ct{s2.txt}\\…\\one id per line};
  \node[act, text width=24mm, minimum height=11mm, double copy shadow={shadow xshift=1.2pt, shadow yshift=-1.2pt}] (sh) at (8.25,1.05) {\textbf{test-without-building}\\-testProductsPath T\\-destination id=UDID\\-only-testing @sK.txt\\-resultBundlePath rK.xcresult};
  \node[act, text width=21mm, minimum height=11mm] (mg) at (11.3,1.05) {\textbf{xcresulttool merge}\\r1.xcresult r2 …\\\dd output-path\\all.xcresult};
  \node[act, text width=20mm, minimum height=11mm] (rd) at (14.15,1.05) {get test-results\\summary · tests\\export attachments};
  \draw[flow] (b1) -- (en); \draw[flow] (en) -- node[lbl, above]{split} (sf); \draw[flow] (sf) -- (sh);
  \draw[flow] (sh) -- node[lbl, above]{×N} (mg); \draw[flow] (mg) -- (rd);
  \node[lbl, anchor=west] at (9.75,1.75) {shards get T, not the project: \ct{-testProductsPath} refuses \ct{-workspace}};
\end{tikzpicture}

\begin{multicols}{2}
\raggedright\footnotesize\setstretch{1.0}

\hd{Build once, test anywhere}
\ct{build-for-testing} writes an \ct{.xctestrun}: an editable plist naming the test bundles, the
host app, the environment and the arguments. Any Mac then runs it with \ct{test-without-building}
— no source, no compiler. Pick tests by path, \ct{Target/Class/method}, with \ct{-only-testing} or
\ct{-skip-testing}; \ct{@file} reads one path per line; given both, \ct{-only-testing} wins.

\hd{Ask the tools before you guess}
\noindent\begin{tabular}{@{}>{\raggedright}p{27mm}@{\hspace{2mm}}l@{}}
which schemes?                & \ct{xcodebuild -list} \\
where can a scheme run?       & \ct{xcodebuild -showdestinations -scheme App} \\
what did a setting become?    & \ct{xcodebuild -showBuildSettings -scheme App} \\
which test plans?             & \ct{xcodebuild -showTestPlans -scheme App} \\
which simulators exist?       & \ct{xcrun simctl list \dd json devices} \\
which binary would run?       & \ct{xcrun \dd find xcodebuild} \\
\end{tabular}
\begin{itemize}
  \trap{No \ct{OS=} in \ct{-destination} means the latest: a new runner image moves every test to a
        new iOS. Pin \ct{OS=}, or the simulator's \ct{id=}.}
\end{itemize}

\hd{Parallel: the simulator you pick never runs a test}
\noindent\begin{tikzpicture}[sheet]
  \node[fil, text width=11mm, align=center, minimum height=7mm] (sim) at (0,0) {chosen\\simulator\\= template};
  \foreach \y/\n/\c in {0.42/1/A, 0/2/B, -0.42/3/C} {
    \node[act, text width=16mm, minimum height=3.4mm, inner sep=0.8pt] (k\n) at (2.2,\y) {clone \n: runner};
    \node[fil, text width=15mm, minimum height=3.4mm, inner sep=0.8pt, align=center] (t\n) at (4.3,\y) {class \c, then next};
    \draw[flow] (sim.east) -- (k\n.west); \draw[flow] (k\n) -- (t\n);
  }
  \node[lbl, anchor=west, align=left] at (5.15,0) {a class per runner:\\class \ct{setUp} once;\\time ≥ slowest class};
\end{tikzpicture}\par
Several \ct{-destination}s each run the \emph{whole} suite; with parallel testing on,
\ct{-parallelize-tests-among-destinations} splits it between them instead.

\hd{Two repetition loops that point opposite ways}
\noindent\begin{tikzpicture}[sheet]
  \node[cmd, anchor=west] at (0,0.42) {-retry-tests-on-failure};
  \foreach \i/\s in {0/ko, 1/ko, 2/ok} { \node[\s] at (2.55+\i*0.28,0.42) {}; }
  \node[lbl, anchor=west, align=left] at (3.35,0.42) {stops at the first \textbf{pass} (≤3): the flake ends \textcolor{sheetGreen!70!black}{\textbf{green}}, hidden};
  \node[cmd, anchor=west] at (0,0) {-run-tests-until-failure};
  \foreach \i in {0,...,3} { \node[ok] at (2.55+\i*0.28,0) {}; }
  \node[lbl] at (3.67,0) {…};
  \node[ko] at (3.92,0) {};
  \node[lbl, anchor=west, align=left] at (4.1,0) {stops at the first \textbf{fail} (≤100): the flake is found};
\end{tikzpicture}\par
\ct{-test-iterations N}: N runs, or the cap on either loop.
\ct{-test-repetition-relaunch-enabled YES}: a fresh process for each run, so no state carries over.

\hd{Environment for the test, through a prefix}
\noindent\begin{tikzpicture}[sheet]
  \node[cmd, anchor=west] (p1) at (0,0.36) {{\color{sheetOrange}TEST\_RUNNER\_}Foo=Bar xcodebuild test};
  \node[cmd, anchor=west] (p2) at (0,0) {{\color{sheetOrange}SIMCTL\_CHILD\_}Foo=Bar simctl launch booted id};
  \draw[hot] (4.5,0.39) -- (5.0,0.39); \draw[hot] (4.5,0.03) -- (5.0,0.03);
  \node[lbl, anchor=west] at (5.05,0.39) {each test runner: \ct{Foo=Bar}};
  \node[lbl, anchor=west] at (5.05,0.03) {the app: \ct{Foo=Bar}};
\end{tikzpicture}\par
The prefix is stripped on the way in. \ct{simctl launch} has no environment flag: the prefix
\emph{is} the interface.

\hd{Read the result bundle, not the log}
\link{\ct{xcresulttool}}{https://keith.github.io/xcode-man-pages/xcresulttool.1.html}
\ct{get test-results summary} (or \ct{tests}) answers in JSON, \ct{help <command>} prints the
schema, and \ct{export attachments} pulls the screenshots out. The exit status only says
\emph{that} it failed: logs show 65 after a failed build or test, 70 for a destination that matches
nothing. Branch on non-zero; ask the bundle why.
\begin{itemize}
  \trap{\ct{xcodebuild … | xcbeautify} without \link{\ct{set -o pipefail}}{https://www.gnu.org/software/bash/manual/html_node/Pipelines.html}:
        the pipeline returns the formatter's status, and a red build reports green.}
\end{itemize}

\hd{Xcode Cloud: your code runs in three scripts}
Only \ct{ci\_post\_clone.sh}, \ct{ci\_pre\_xcodebuild.sh} and \ct{ci\_post\_xcodebuild.sh}, from
\link{\ct{ci\_scripts/}}{https://developer.apple.com/documentation/xcode/writing-custom-build-scripts}
beside the project. The last runs \emph{even when the build failed}: check before it uploads or
notifies. No \ct{sudo}; secrets print as asterisks.

\columnbreak

\hd{Export: \ct{method} answers ``who may install this?''}
\noindent\begin{tikzpicture}[sheet]
  \node[st, text width=17mm, minimum height=9mm] (m1) at (0.95,0) {\ct{\textbf{debugging}}\\(the default)\\the team's devices\\in the profile};
  \node[st, text width=21mm, minimum height=11mm] (m2) at (3.55,0) {\ct{\textbf{release-testing}}\\(ad hoc) only the devices\\listed in the profile:\\\textbf{≤100} per family a year};
  \node[st, text width=21mm, minimum height=13mm, draw=sheetGreen!70!black, fill=sheetGreen!10] (m3) at (6.4,0) {\ct{\textbf{app-store-connect}}\\TestFlight first,\\then anyone\\after review};
  \draw[flow] (m1) -- (m2); \draw[flow] (m2) -- (m3);
  \node[lbl, anchor=north west, align=left] at (-0.05,-0.72) {off the ladder: \ct{enterprise} (your own staff, no store) · \ct{developer-id} (a Mac app outside the store)};
\end{tikzpicture}\par
The 100 is per product family per membership year; \link{disabling a device does not give the slot
back}{https://developer.apple.com/help/account/devices/devices-overview/}, only renewal does.
Manual signing needs \ct{provisioningProfiles}: every bundle id mapped to its profile.
\begin{itemize}
  \trap{\ct{manageAppVersionAndBuildNumber} raises a used build number on upload, in a staged copy,
        never in the archive. Turn it off when CI sets the number.}
\end{itemize}

\hd{Signing with nobody at the keyboard}
\ct{-allowProvisioningUpdates} lets \ct{xcodebuild} talk to the developer website: for automatic
signing it creates certificates, app IDs and profiles; for manual, it downloads the profiles. It
needs an account signed in to Xcode, or an App Store Connect API key in its place:
\ct{-authenticationKeyPath}, \ct{-authenticationKeyID}, \ct{-authenticationKeyIssuerID}. The same
key uploads.

\hd{Swift packages, plugins, macros on CI}
\begin{itemize}
  \item \ct{Package.resolved} is the lockfile: commit it;
        \link{\ct{-disableAutomaticPackageResolution}}{https://developer.apple.com/documentation/xcode/building-swift-packages-or-apps-that-use-them-in-continuous-integration-workflows}
        makes CI build exactly those versions.
  \item Checkouts land in \ct{DerivedData/SourcePackages}; \ct{-clonedSourcePackagesDirPath P} puts
        them where a cache survives.
  \item \ct{-scmProvider system}: the Mac's own git, its configuration and SSH keys; the host must be
        in \ct{known\_hosts}.
  \item Nobody is there to click ``Trust \& Enable'': \ct{-skipPackagePluginValidation} and
        \ct{-skipMacroValidation} trust every plugin and macro unseen.
\end{itemize}

\hd{A simulator is a userspace, not a virtual machine}
\link{The iOS userspace}{https://developer.apple.com/videos/play/wwdc2019/418/} — its own
\ct{launchd}, daemons and frameworks — runs on the \emph{Mac's} kernel and filesystem, with a home
of its own. CPU and memory limits are not simulated: on a Mac with a terabyte of RAM the
``iPhone'' can allocate a terabyte. Simulator timings say nothing about a phone.\par\vspace{2pt}
\noindent\begin{tikzpicture}[sheet]
  \node[fil, text width=13mm, align=center] (cr) at (0.75,0) {\ct{create}\\prints a UDID};
  \node[st, text width=13mm, minimum height=6mm] (sd) at (3.0,0) {\textbf{Shutdown}};
  \node[st, text width=13mm, minimum height=6mm, draw=sheetGreen!70!black, fill=sheetGreen!10] (bt) at (5.6,0) {\textbf{Booted}};
  \draw[flow] (cr) -- (sd);
  \draw[flow] ([yshift=1mm]sd.east) -- node[lbl, above]{\ct{boot}} ([yshift=1mm]bt.west);
  \draw[flow] ([yshift=-1mm]bt.west) -- node[lbl, below]{\ct{shutdown}} ([yshift=-1mm]sd.east);
  \node[lbl, anchor=north, align=center] at (3.0,-0.38) {\ct{erase} · \ct{clone}: only here;\\clone = APFS copy-on-write};
  \node[lbl, anchor=west, align=left] at (6.45,0) {\ct{install} · \ct{launch}\\\ct{openurl} · \ct{push}\\\ct{privacy} · \ct{io}};
\end{tikzpicture}\par
A script can do to a simulator what a phone never allows: \ct{privacy … grant} (no alert to tap),
\ct{push} an \ct{.apns} file, \ct{keychain … add-root-cert} (trust your proxy),
\ct{status\_bar … override \dd time 9:41} (one clock in every screenshot). Script with UDIDs;
\ct{booted} means the one booted device. Phones: \ct{devicectl}, \ct{\dd device} by name or UDID;
parse its \ct{\dd json-output}, never its text.

\hd{Which Xcode answers}
\noindent\begin{tikzpicture}[sheet]
  \node[fil, text width=19mm, align=center] (sh) at (1.03,0) {\ct{/usr/bin/xcodebuild}\\a shim};
  \node[st, text width=28mm] (dv) at (4.05,0.3) {\ct{DEVELOPER\_DIR=…/Xcode\_27.app}\\this process only};
  \node[st, text width=28mm] (xs) at (4.05,-0.3) {\ct{sudo xcode-select -s …}\\every user, until changed};
  \node[fil, text width=16.5mm, align=center] (tl) at (7.0,0) {\ct{…/Contents/}\\\ct{Developer/usr/bin}\\the real tool};
  \draw[flow] (sh.east) -- node[lbl, anchor=south east, pos=0.5]{set?} (dv.west);
  \draw[flow] (sh.east) -- node[lbl, anchor=north east, pos=0.5]{else} (xs.west);
  \draw[flow] (dv.east) -- (tl.west); \draw[flow] (xs.east) -- (tl.west);
\end{tikzpicture}\par
Pin the version: \ct{xcodes} installs the one \ct{.xcode-version} names. A fresh Mac:
\ct{-runFirstLaunch}; download a platform once with \ct{-downloadPlatform iOS -exportPath D}, then
\ct{-importPlatform} it on every other Mac.

\begin{htmlonly}
\section{Further reading}
\begin{itemize}
  \item \href{https://keith.github.io/xcode-man-pages/}{Xcode man pages} — every man page Xcode ships, browsable: xcodebuild, xcresulttool, xcrun, xcode-select and the rest
  \item \href{https://developer.apple.com/videos/play/wwdc2019/418/}{Getting the Most Out of Simulator (WWDC19)} — the simulator's insides from the people who build it: a userspace on the Mac kernel, why limits are not simulated, clones as APFS copies
  \item \href{https://asciiwwdc.com/2018/sessions/403}{What's New in Testing (WWDC18, transcript)} — how parallel testing works: a template simulator, named clones, classes handed out one at a time
  \item \href{https://developer.apple.com/documentation/xcode/building-swift-packages-or-apps-that-use-them-in-continuous-integration-workflows}{Swift packages in continuous integration} — Apple's own checklist: the lockfile, SSH access, the Mac's git, private package hosts
  \item \href{https://saagarjha.com/blog/2019/01/11/mocking-capabilities-in-the-ios-simulator/}{Mocking Capabilities in the iOS Simulator} — \texttt{SIMCTL\_CHILD\_} taken to its limit: a capabilities plist that changes what the simulated device claims it can do
  \item \href{https://nshipster.com/simctl/}{simctl (NSHipster)} — a tour of the subcommands, from pasteboard syncing to switching a simulator's language without opening Settings
\end{itemize}
\end{htmlonly}

\end{multicols}

\noindent{\footnotesize\color{sheetGrey}\textit{Related:} \texttt{ci-pipeline-design} ·
\texttt{fastlane-deep} · \texttt{github-actions-deep} · Code signing \& iOS CI/CD ·
\texttt{macos-keychain-security-cli} · Flaky tests · CI · coverage · TDD ·
\texttt{xctest-and-swift-testing} · Crashes \& symbolication · Modularization \& Swift Package
Manager · \texttt{macos-power-user-cli}}

\end{document}
